% TOP_PROBLEM: {"id": 1195, "title": "Finiteness of Finitely Presented Periodic Groups", "category_id": 4, "category": {"id": 4, "name": "algebra", "display_name": "Algebra", "description": "Group theory, ring theory, field theory, and algebraic structures.", "slug": "algebra", "order_index": 4, "created_at": "2026-07-31T15:26:25.671Z"}, "status": "open"}
% FIRST_POSTED: 2026-09-27
% !TeX program = lualatex
% ATTEMPT_STATUS: unresolved
\documentclass[11pt]{article}
\usepackage[margin=25mm]{geometry}
\usepackage{fontspec}
\setmainfont{Latin Modern Roman}
\usepackage{amsmath,amssymb,mathtools}
\usepackage{unicode-math}
\setmathfont{Latin Modern Math}
\usepackage{xurl,hyperref}
\hypersetup{colorlinks=true,urlcolor=blue}
\setcounter{secnumdepth}{0}
\setlength{\parindent}{0pt}
\setlength{\parskip}{5pt}
\setlength{\emergencystretch}{3em}
\begin{document}
\section*{\texorpdfstring{Finiteness of Finitely Presented Periodic Groups}{Finiteness of Finitely Presented Periodic Groups}}\label{1195}
\subsection{Definitions and mathematical statement}
A group is finitely presented if it has a presentation $G=\langle s_1,\ldots,s_d\mid r_1,\ldots,r_m\rangle$ with finitely many generators and relators. It is periodic, or torsion, if
\[
 \forall g\in G\ \exists n=n(g)\ge1:\ g^n=1.
\]
Does every finitely presented periodic group have finitely many elements? Equivalently, does an infinite finitely presented torsion group exist? There is no single exponent required to kill every element: demanding a uniform $n$ would impose a stronger condition. Finite generation alone is also weaker than finite presentability, so a construction specified by infinitely many relations does not answer this question.
\par\smallskip\textit{Formulation and concept references:} \href{https://groupprops.subwiki.org/wiki/Finitely_presented_periodic_group}{[S1]}.

\subsection{Short English statement}
Can an infinite group have both a finite presentation and finite order for every one of its elements?
\subsection{Research attempt: finite presentation versus successive torsion relations}
\textbf{Status and the indispensable hypothesis.} This attempt does not prove that every finitely presented periodic group is finite, and it does not construct an infinite finitely presented periodic group. Infinite finitely generated torsion groups are not counterexamples. A current primary example is the work of Kionke and Schesler [E1], which constructs finitely generated residually finite infinite torsion groups with further properties; its abstract does not replace finite generation by finite presentation. The primary discussion of Haettel and Osajda [E2] also treats torsion groups in geometric settings. These sources were checked on 27 September 2026 without finding a proof of the selected general assertion.

Here ``periodic'' means $\forall g\ \exists n(g)\geq1$ with $g^{n(g)}=1$. There is no single common exponent in the hypothesis. The work below proves finiteness for a substantial elementary subclass and isolates exactly why adding torsion relations indefinitely usually does not produce a finite presentation.

\textbf{Lemma 319.1 (finite generation in a finite-index subgroup).} If $G$ is generated by a finite symmetric set $S$ and $H\leq G$ has finite index, then $H$ is finitely generated.

\emph{Proof.} Choose representatives $T$ for the right cosets $H\backslash G$, including $1$. For $t\in T$ and $s\in S$, let $\overline{ts}\in T$ represent $Hts$. Then
\[
 t s(\overline{ts})^{-1}\in H.
\]
These finitely many elements generate $H$. For a word $s_1\cdots s_r\in H$, insert the representatives of its successive prefix cosets between consecutive letters. The resulting product telescopes into elements of the displayed form, and the final representative is $1$. $\square$

\textbf{Theorem 319.2 (the virtually solvable subcase).} A finitely generated periodic virtually solvable group is finite.

\emph{Proof.} First suppose $G$ is solvable. Induct on its derived length. If it is abelian and generated by $g_1,\ldots,g_d$, of respective finite orders $n_1,\ldots,n_d$, every element has a representative
\[
 g_1^{a_1}\cdots g_d^{a_d},\qquad 0\leq a_i<n_i.
\]
Thus $|G|\leq\prod_i n_i$. For larger derived length, the abelianization $G/G'$ is finitely generated and periodic, hence finite by this base case. Therefore $G'$ has finite index and is finitely generated by Lemma 319.1. It remains periodic and has smaller derived length, so induction makes it finite. A finite extension of a finite group is finite.

If $G$ is only virtually solvable, choose a solvable subgroup $H$ of finite index. Lemma 319.1 makes it finitely generated, and periodicity passes to subgroups. The preceding argument makes $H$ finite, and then $G$ finite. $\square$

No finite-presentation assumption was needed in this proof. Consequently any counterexample to the selected assertion would necessarily lie outside the virtually solvable class. The finite-index step in the induction is essential: a subgroup of a finitely generated group need not generally be finitely generated.

\textbf{Two elementary warnings about orders.} Knowing that each member of a finite generating set has finite order does not even show that the group is periodic. For example, the transformations
\[
 a(x)=-x,\qquad b(x)=2-x
\]
of the real line are involutions, but $ab(x)=x-2$ has infinite order. They exhibit the obstruction already in a quotient of the presentation $\langle a,b\mid a^2=b^2=1\rangle$.

By contrast, if every element of a group has order dividing two, then
\[
 ab=(ab)^{-1}=b^{-1}a^{-1}=ba.
\]
Such a group is abelian, and a $d$-generated example has at most $2^d$ elements. This proof uses the order condition on the product $ab$, not just on the generators. It illustrates both a solved bounded-exponent subcase and the danger of replacing the quantifier over all words by a finite list of obvious words.

\textbf{Lemma 319.3 (finite presentation gives stabilization on fixed generators).} Let $F$ be a finitely generated free group and let
\[
 N_1\subseteq N_2\subseteq\cdots\lhd F,
 \qquad N=\bigcup_iN_i.
\]
If $F/N$ is finitely presented, then $N=N_j$ for some $j$.

\emph{Proof.} For a fixed finite generating set of a finitely presented group, the kernel of the map from the corresponding free group is finitely normally generated. To justify the qualification about generators, start with any finite presentation, express its generators as words in the chosen generators and express the chosen generators as words in the original ones. Adjoin the chosen generators and these finitely many relations to the presentation, then eliminate the original generators by substitution. This produces a finite presentation on the chosen set.

Apply this observation to $F/N$, writing $N$ as the normal closure of finitely many words $r_1,\ldots,r_k$. Each $r_i$ lies in some $N_{j_i}$. For $j=\max_i j_i$, normality gives $N\subseteq N_j$, while the reverse inclusion holds by definition. $\square$

In particular, a strictly increasing sequence of normal subgroups on a fixed finite generating set cannot have a finitely presented direct-limit quotient. This is a useful test for an alleged construction: if genuinely new independent relations are being added forever, finite presentation at the end has not been achieved.

\textbf{What the torsion-by-enumeration construction does and does not prove.} Enumerate the words of $F$ as $w_1,w_2,\ldots$. One can try to construct successive quotients by imposing
\[
 w_i^{n_i}=1
\]
after stage $i$, choosing the exponents to preserve desired features. In the direct limit every group element is represented by some $w_i$, hence has finite order. This part of the argument is valid.

There are then two independent obligations. One must prove the limit is infinite, and one must prove it is finitely presented. A useful observation about the first obligation is that if every stage quotient $F/N_i$ is infinite, the limit quotient cannot be finite. Indeed, every finite group has a finite presentation: use a generator for each group element and its finite multiplication table as relations. If the limit were finite, Lemma 319.3 would make the sequence stabilize at that finite quotient, contradicting the assumed infinitude of every stage.

However, ensuring an infinite stage quotient after each added relation is a real mathematical task, not a formal consequence of choosing a large exponent. And if the construction also proves strict growth $N_i\ne N_{i+1}$ at arbitrarily late stages, Lemma 319.3 shows that its infinite periodic limit is \emph{not} finitely presented. This explains why a successful infinite-torsion construction may still miss exactly the hypothesis in the selected problem.

The lemma does not prove the conjecture by itself. A finitely presented periodic group already has all its torsion identities as consequences of finitely many relators. Listing those consequences one at a time need not enlarge the normal subgroup at all. The missing argument would have to rule out this possibility for an infinite finitely presented group, which is precisely what has not been established.

\textbf{Finite quotients and residual finiteness.} Every finite quotient of a periodic group is periodic, but the converse implication from finite quotients is false. For instance, all finite quotients of $\mathbb Z$ are finite cyclic groups, while $1\in\mathbb Z$ has infinite order. Even residual finiteness does not repair that inference: $\mathbb Z$ is residually finite. A bound on the orders of all finite images of a fixed element would be extra information, and a uniform exponent across all elements would be stronger still.

The primary constructions in [E1] also show why residual finiteness should not be advertised as a substitute for finite presentation in an argument for finiteness: they retain residual finiteness while producing infinite torsion groups. The present attempt does not use an unproved claim that such examples are finitely presented or that their additional invariants force finite presentation.

\textbf{A finite algebraic diagnostic.} The accompanying exact script checks the affine involution example using integer affine maps, verifies the exponent-two counting argument on small elementary abelian groups, and enumerates residue representatives for selected finite abelian presentations. The latter calculations test the base case and presentations of finite quotients only. They cannot certify periodicity of every word in an arbitrary infinite presentation, and no finite list of observed finite orders is presented as doing so.

\textbf{Outcome and first gap.} Periodic finitely generated virtually solvable groups are finite, and a fixed-generator quotient obtained through infinitely many strict relation enlargements cannot be finitely presented. These results rule out several natural shortcuts on both sides of the question. The unresolved possibility is an infinite non-virtually-solvable group whose torsion identities all follow from a finite presentation. Neither finite quotient calculations nor the existence of finitely generated torsion examples determines whether such a group exists. The selected assertion remains unresolved by this attempt.

\subsection{Sources}
\begin{itemize}
\item[S1]Groupprops contributors, Finitely presented periodic group, Subwiki.
\url{https://groupprops.subwiki.org/wiki/Finitely_presented_periodic_group}.
\textit{Formulation source.}
\item[E1]Steffen Kionke and Eduard Schesler, \emph{Hereditarily just-infinite torsion groups with positive first $\ell^2$-Betti number}, arXiv:2401.04542 (2024).
\url{https://arxiv.org/abs/2401.04542}.
\textit{Primary source consulted 27 September 2026; its examples are stated as finitely generated, not as finitely presented.}
\item[E2]Thomas Haettel and Damian Osajda, \emph{Locally elliptic actions, torsion groups, and nonpositively curved spaces}, arXiv:2110.12431.
\url{https://arxiv.org/abs/2110.12431}.
\textit{Primary source consulted 27 September 2026 for the scope of geometric restrictions on torsion groups.}
\end{itemize}
\end{document}
